\chapter{1989年全国卷（理）}

\section{单选题}

\begin{problem}
如果 \(I = \{a, b, c, d, e\}\)，\(M = \{a, c, d\}\)，\(N = \{b, d, e\}\)，其中 \(I\) 是全集，那么 \(\overline{M} \cap \overline{N}\) 等于（\quad）

\choices
    {\(\varnothing\)}
    {\(\{d\}\)}
    {\(\{a, c\}\)}
    {\(\{b, e\}\)}
\end{problem}

\begin{answer}
A
\end{answer}

\begin{solution}
由于 \(I\) 是全集，\(I\) 中除去 \(M\) 的元素为 \(b,e\)，除去 \(N\) 的元素为 \(a,c\)，所以 \(\overline{M}=\{b,e\}\)，\(\overline{N}=\{a,c\}\). 因而 \(\overline{M}\cap\overline{N}=\varnothing\)，故选 A.
\end{solution}

\begin{problem}
与函数 \(y = x\) 有相同图象的一个函数是（\quad）

\choices
    {\(y = \sqrt{x^2}\)}
    {\(y = \frac{x^2}{x}\)}
    {\(y = a^{\log_a x}\)（其中 \(a > 0\)，\(a \neq 1\)）}
    {\(y = \log_a a^x\)（其中 \(a > 0\)，\(a \neq 1\)）}
\end{problem}

\begin{answer}
D
\end{answer}

\begin{solution}
选项 A 中 \(y=\sqrt{x^2}=|x|\)，当 \(x<0\) 时不等于 \(x\)；选项 B 中 \(y=\frac{x^2}{x}=x\) 只在 \(x\neq0\) 时有定义，定义域不是 \(\R\)；选项 C 中 \(y=a^{\log_a x}=x\) 的定义域为 \(x>0\)，也不是 \(\R\). 选项 D 中，由于 \(a>0\) 且 \(a\neq1\)，对任意 \(x\in\R\) 都有 \(y=\log_a a^x=x\)，故选 D.
\end{solution}

\begin{problem}
如果圆锥的底面半径为 \(\sqrt{2}\)，高为 \(2\)，那么它的侧面积是（\quad）

\choices
    {\(4\sqrt{3}\pi\)}
    {\(2\sqrt{2}\pi\)}
    {\(2\sqrt{3}\pi\)}
    {\(4\sqrt{2}\pi\)}
\end{problem}

\begin{answer}
C
\end{answer}

\begin{solution}
设圆锥的母线长为 \(l\)，则由勾股定理，\(l=\sqrt{(\sqrt{2})^2+2^2}=\sqrt{6}\). 所以圆锥的侧面积为 \(\pi\sqrt{2}\cdot\sqrt{6}=2\sqrt{3}\pi\)，故选 C.
\end{solution}

\begin{problem}
\(\cos \left[\arcsin \left(-\frac{4}{5}\right) - \arccos \left(-\frac{3}{5}\right)\right]\) 的值等于（\quad）

\choices
    {\(-1\)}
    {\(-\frac{7}{25}\)}
    {\(\frac{7}{25}\)}
    {\(-\frac{\sqrt{10}}{5}\)}
\end{problem}

\begin{answer}
A
\end{answer}

\begin{solution}
设 \(\alpha=\arcsin\frac{4}{5}\)，则 \(\alpha\in(0,\frac{\pi}{2})\)，且 \(\cos\alpha=\frac{3}{5}\). 因而 \(\arcsin(-\frac45)=-\alpha\)，\(\arccos(-\frac35)=\pi-\alpha\). 所以
\[
\cos\left[\arcsin\left(-\frac45\right)-\arccos\left(-\frac35\right)\right]=\cos(-\pi)=-1,
\]
故选 A.
\end{solution}

\begin{problem}
已知 \(\{a_n\}\) 是等比数列，如果 \(a_1 + a_2 + a_3 = 18\)，\(a_2 + a_3 + a_4 = -9\)，且 \(S_n = a_1 + a_2 + \cdots + a_n\)，那么 \(\displaystyle\lim_{n \to \infty} S_n\) 的值等于（\quad）

\choices
    {\(8\)}
    {\(16\)}
    {\(32\)}
    {\(48\)}
\end{problem}

\begin{answer}
B
\end{answer}

\begin{solution}
设等比数列的公比为 \(q\). 因为 \(a_2+a_3+a_4=q(a_1+a_2+a_3)\)，所以 \(q=\frac{-9}{18}=-\frac12\). 又
\[
18=a_1(1+q+q^2)=a_1\left(1-\frac12+\frac14\right)=\frac34a_1,
\]
从而 \(a_1=24\). 由于 \(|q|=\frac12<1\)，无穷等比级数收敛，故
\[
\lim_{n\to\infty}S_n=\frac{a_1}{1-q}=\frac{24}{1+\frac12}=16,
\]
故选 B.
\end{solution}

\begin{problem}
如果 \(\left|\cos \theta\right| = \frac{1}{5}\)，\(\frac{5}{2}\pi < \theta < 3\pi\)，那么 \(\sin \frac{\theta}{2}\) 的值等于（\quad）

\choices
    {\(-\frac{\sqrt{10}}{5}\)}
    {\(\frac{\sqrt{10}}{5}\)}
    {\(-\frac{\sqrt{15}}{5}\)}
    {\(\frac{\sqrt{15}}{5}\)}
\end{problem}

\begin{answer}
C
\end{answer}

\begin{solution}
由 \(\frac{5\pi}{2}<\theta<3\pi\)，可知 \(\frac{\pi}{2}<\theta-2\pi<\pi\)，因此 \(\cos\theta<0\). 结合 \(|\cos\theta|=\frac15\)，得 \(\cos\theta=-\frac15\). 又 \(\frac{5\pi}{4}<\frac{\theta}{2}<\frac{3\pi}{2}\)，所以 \(\sin\frac{\theta}{2}<0\). 由半角公式，
\[
\sin^2\frac{\theta}{2}=\frac{1-\cos\theta}{2}=\frac{1+\frac15}{2}=\frac35,
\]
故 \(\sin\frac{\theta}{2}=-\frac{\sqrt{15}}5\)，选 C.
\end{solution}

\begin{problem}
设复数 \(z\) 满足关系式 \(z + |z| = 2 + \i\)，那么 \(z\) 等于（\quad）

\choices
    {\(-\frac{3}{4} + \i\)}
    {\(\frac{3}{4} - \i\)}
    {\(-\frac{3}{4} - \i\)}
    {\(\frac{3}{4} + \i\)}
\end{problem}

\begin{answer}
D
\end{answer}

\begin{solution}
设 \(z=x+y\i\)，其中 \(x,y\in\R\). 比较实部和虚部，得 \(x+|z|=2\)，\(y=1\). 令 \(r=|z|\)，则 \(r=2-x\)，且 \(r^2=x^2+1\). 因而
\[
(2-x)^2=x^2+1,
\]
解得 \(x=\frac34\)，再由 \(r=2-x\) 得 \(r=\frac54\). 所以 \(z=\frac34+\i\)，故选 D.
\end{solution}

\begin{problem}
已知球的两个平行截面的面积分别为 \(5\pi\) 和 \(8\pi\)，它们位于球心的同一侧，且相距为 \(1\)，那么这个球的半径是（\quad）

\choices
    {\(4\)}
    {\(3\)}
    {\(2\)}
    {\(5\)}
\end{problem}

\begin{answer}
B
\end{answer}

\begin{solution}
设球的半径为 \(R\)，球心到面积为 \(5\pi\) 和 \(8\pi\) 的截面的距离分别为 \(d_5\) 和 \(d_8\). 由截面圆面积公式，
\[
d_5^2=R^2-5,\qquad d_8^2=R^2-8.
\]
面积为 \(8\pi\) 的截面更靠近球心，且两截面在球心同一侧，所以 \(d_5-d_8=1\). 因此
\[
\sqrt{R^2-5}-\sqrt{R^2-8}=1.
\]
令 \(u=\sqrt{R^2-8}\)，则 \(\sqrt{u^2+3}=u+1\). 两边平方得 \(u^2+3=u^2+2u+1\)，从而 \(u=1\)，所以 \(R^2=9\)，\(R=3\)，故选 B.
\end{solution}

\begin{problem}
已知椭圆的极坐标方程是 \(\rho = \frac{5}{3 - 2\cos \theta}\)，那么它的短轴长是（\quad）

\choices
    {\(\frac{10}{3}\)}
    {\(\sqrt{5}\)}
    {\(2\sqrt{5}\)}
    {\(2\sqrt{3}\)}
\end{problem}

\begin{answer}
C
\end{answer}

\begin{solution}
由 \(x=\rho\cos\theta\) 及题设方程，得 \(3\rho-2x=5\)，即 \(3\rho=5+2x\). 两边平方并利用 \(\rho^2=x^2+y^2\)，有
\[
9(x^2+y^2)=(5+2x)^2,
\]
整理得 \(5x^2+9y^2-20x-25=0\). 配方可得
\[
5(x-2)^2+9y^2=45,
\]
即
\[
\frac{(x-2)^2}{9}+\frac{y^2}{5}=1.
\]
所以椭圆的长半轴为 \(3\)，短半轴为 \(\sqrt5\)，短轴长为 \(2\sqrt5\)，故选 C.
\end{solution}

\begin{problem}
如果双曲线 \(\frac{x^2}{64} - \frac{y^2}{36} = 1\) 上一点 \(P\) 到它的右焦点的距离是 \(8\)，那么点 \(P\) 到它的右准线的距离是（\quad）

\choices
    {\(10\)}
    {\(\frac{32\sqrt{7}}{7}\)}
    {\(2\sqrt{7}\)}
    {\(\frac{32}{5}\)}
\end{problem}

\begin{answer}
D
\end{answer}

\begin{solution}
双曲线的实半轴长为 \(a=8\)，虚半轴长为 \(b=6\)，所以焦半距
\[
c=\sqrt{a^2+b^2}=\sqrt{64+36}=10,
\]
离心率为 \(e=\frac{c}{a}=\frac54\). 右准线为 \(x=\frac{a}{e}\). 题设中的点 \(P\) 必在右支上，因为左支上的点到右焦点的距离不小于 \(18\)，不可能等于 \(8\). 由双曲线的焦点—准线定义，\(PF=e\cdot d\)，其中 \(d\) 是 \(P\) 到右准线的距离. 因而
\[
d=\frac{PF}{e}=\frac{8}{\frac54}=\frac{32}{5},
\]
故选 D.
\end{solution}

\begin{problem}
已知 \(f(x) = 8 + 2x - x^2\)，如果 \(g(x) = f(2 - x^2)\)，那么 \(g(x)\)（\quad）

\choices
    {在区间 \((-1,0)\) 上是减函数}
    {在区间 \((0,1)\) 上是减函数}
    {在区间 \((-2,0)\) 上是增函数}
    {在区间 \((0,2)\) 上是增函数}
\end{problem}

\begin{answer}
A
\end{answer}

\begin{solution}
先化简复合函数：
\[
g(x)=8+2(2-x^2)-(2-x^2)^2=8+2x^2-x^4.
\]
所以 \(g'(x)=4x(1-x^2)\). 在 \((-1,0)\) 上，\(x<0\) 且 \(1-x^2>0\)，故 \(g'(x)<0\)，函数为减函数. 在 \((0,1)\) 上函数为增函数；在 \((-2,0)\) 和 \((0,2)\) 上，导数都发生变号，不能分别判断为增函数. 因此只有选项 A 正确.
\end{solution}

\begin{problem}
由数字 \(1\)，\(2\)，\(3\)，\(4\)，\(5\) 组成没有重复数字的五位数，其中小于 \(50000\) 的偶数共有（\quad）

\choices
    {\(60\text{ 个}\)}
    {\(48\text{ 个}\)}
    {\(36\text{ 个}\)}
    {\(24\text{ 个}\)}
\end{problem}

\begin{answer}
C
\end{answer}

\begin{solution}
五位数小于 \(50000\)，首位只能取 \(1,2,3,4\)；末位为偶数，末位只能取 \(2,4\). 若首位取 \(1\) 或 \(3\)，有 \(2\) 种首位选择和 \(2\) 种末位选择，中间三位从剩下的三个数字中排列，有 \(3!\) 种，共有 \(2\cdot2\cdot3!=24\) 个. 若首位取 \(2\) 或 \(4\)，末位只能取另一个偶数，有 \(2\) 种首位选择和 \(1\) 种末位选择，中间三位仍有 \(3!\) 种，共有 \(2\cdot1\cdot3!=12\) 个. 所以总数为 \(24+12=36\)，故选 C.
\end{solution}

\section{填空题}

\begin{problem}
方程 \(\sin x - \sqrt{3}\cos x = \sqrt{2}\) 的解集是\fillinblank{}.
\end{problem}

\begin{answer}
\(x = \frac{7\pi}{12} + 2k\pi\) 或 \(x = \frac{13\pi}{12} + 2k\pi\)，\(k \in \Z\)
\end{answer}

\begin{solution}
利用正弦差角公式，
\[
\sin x-\sqrt3\cos x=2\sin\left(x-\frac{\pi}{3}\right).
\]
所以原方程等价于 \(\sin(x-\frac{\pi}{3})=\frac{\sqrt2}{2}\). 因而
\[
x-\frac{\pi}{3}=\frac{\pi}{4}+2k\pi\quad\text{或}\quad x-\frac{\pi}{3}=\frac{3\pi}{4}+2k\pi,\qquad k\in\Z.
\]
解得 \(x=\frac{7\pi}{12}+2k\pi\) 或 \(x=\frac{13\pi}{12}+2k\pi\)，其中 \(k\in\Z\).
\end{solution}

\begin{problem}
不等式 \(|x^2 - 3x| > 4\) 的解集是\fillinblank{}.
\end{problem}

\begin{answer}
\((-\infty,-1) \cup (4,+\infty)\)
\end{answer}

\begin{solution}
由绝对值不等式的定义，原不等式等价于
\[
x^2-3x>4\quad\text{或}\quad x^2-3x<-4.
\]
第一种情况给出 \(x^2-3x-4>0\)，即 \((x-4)(x+1)>0\)，所以 \(x<-1\) 或 \(x>4\). 第二种情况给出 \(x^2-3x+4<0\)，其判别式为 \(9-16=-7<0\)，且二次项系数为正，故无实数解. 因此解集为 \((-\infty,-1)\cup(4,+\infty)\).
\end{solution}

\begin{problem}
函数 \(y = \frac{\e^x - 1}{\e^x + 1}\) 的反函数的定义域是\fillinblank{}.
\end{problem}

\begin{answer}
\((-1,1)\)
\end{answer}

\begin{solution}
令 \(t=\e^x\)，则 \(t>0\)，且 \(y=\frac{t-1}{t+1}\). 由 \(t>0\) 可知 \(-1<\frac{t-1}{t+1}<1\). 反过来，对任意 \(y\in(-1,1)\)，由
\[
y=\frac{t-1}{t+1}
\]
可得 \(t=\frac{1+y}{1-y}>0\)，从而存在 \(x=\ln\frac{1+y}{1-y}\). 所以原函数的值域为 \((-1,1)\)，即反函数的定义域为 \((-1,1)\).
\end{solution}

\begin{problem}
已知 \((1 - 2x)^7 = a_0 + a_1x + a_2x^2 + \cdots + a_7x^7\)，那么 \(a_1 + a_2 + \cdots + a_7 = \)\fillinblank{}.
\end{problem}

\begin{answer}
\(-2\)
\end{answer}

\begin{solution}
令 \(x=1\)，题设展开式给出
\[
a_0+a_1+\cdots+a_7=(1-2)^7=-1.
\]
再令 \(x=0\)，得 \(a_0=(1-0)^7=1\). 因此
\[
a_1+a_2+\cdots+a_7=-1-a_0=-2.
\]
\end{solution}

\begin{problem}
已知 \(A\) 和 \(B\) 是两个命题，如果 \(A\) 是 \(B\) 的充分条件，那么 \(B\) 是 \(A\) 的\fillinblank{}条件；\(\overline{A}\) 是 \(\overline{B}\) 的\fillinblank{}条件.
\end{problem}

\begin{answer}
必要；必要
\end{answer}

\begin{solution}
“\(A\) 是 \(B\) 的充分条件”表示 \(A\Rightarrow B\)，所以 \(B\) 是 \(A\) 的必要条件. 对这个蕴含关系取逆否命题，得 \(\overline{B}\Rightarrow\overline{A}\)，所以 \(\overline{A}\) 是 \(\overline{B}\) 的必要条件. 两个空均填“必要”.
\end{solution}

\begin{problem}
如图，已知圆柱的底面半径是 \(3\)，高是 \(4\)，\(A\)，\(B\) 两点分别在两底面的圆周上，并且 \(AB = 5\)，那么直线 \(AB\) 与轴 \(OO'\) 之间的距离等于\fillinblank{}.

\bitmapinlinefigure[0.34\linewidth][max width=0.34\linewidth]{img/1989/national_science/q18_fig1.png}
\end{problem}

\begin{answer}
\(\frac{3\sqrt{3}}{2}\)
\end{answer}

\begin{solution}
把 \(A\)，\(B\) 沿轴方向投影到同一个水平面，记所得点为 \(A_0\)，\(B_0\)，并记轴 \(OO'\) 在该水平面上的投影点为 \(O_0\). 则 \(O_0A_0=O_0B_0=3\). 由于两底面间的高为 \(4\)，有
\[
A_0B_0^2=AB^2-4^2=5^2-4^2=9,
\]
所以 \(A_0B_0=3\). 在等腰三角形 \(O_0A_0B_0\) 中，腰长为 \(3\)，底边长为 \(3\)，从 \(O_0\) 到直线 \(A_0B_0\) 的距离为
\[
\sqrt{3^2-\left(\frac32\right)^2}=\frac{3\sqrt3}{2}.
\]
直线 \(AB\) 与轴 \(OO'\) 的最短距离正好等于其水平投影直线 \(A_0B_0\) 到 \(O_0\) 的距离，因此所求距离为 \(\frac{3\sqrt3}{2}\).
\end{solution}

\section{解答题}

\begin{problem}
证明：\(\tan \frac{3x}{2} - \tan \frac{x}{2} = \frac{2\sin x}{\cos x + \cos 2x}\).
\end{problem}

\begin{solution}
在两边均有意义的范围内，利用正切差角公式和积化和差公式，有
\[
\tan \frac{3x}{2} - \tan \frac{x}{2} = \frac{\sin\left(\frac{3x}{2}-\frac{x}{2}\right)}{\cos\frac{3x}{2}\cos\frac{x}{2}} = \frac{\sin x}{\dfrac{1}{2}\left(\cos 2x+\cos x\right)} = \frac{2\sin x}{\cos x+\cos 2x}.
\]
所以等式成立.
\end{solution}

\begin{problem}
如图，在平行六面体 \(ABCD - A_1B_1C_1D_1\) 中，已知 \(AB = 5\)，\(AD = 4\)，\(AA_1 = 3\)，\(AB \perp AD\)，\(\angle A_1AB = \angle A_1AD = \frac{\pi}{3}\).
\begin{enumerate}
\item 求证：顶点 \(A_1\) 在底面 \(ABCD\) 的射影 \(O\) 在 \(\angle BAD\) 的平分线上；
\item 求这个平行六面体的体积.
\end{enumerate}

\bitmapinlinefigure[0.50\linewidth][max width=0.50\linewidth]{img/1989/national_science/q20_fig1.png}
\end{problem}

\begin{solution}
\begin{enumerate}
\item 设 \(O\) 为 \(A_1\) 在底面 \(ABCD\) 上的射影，在底面内以 \(AB\)，\(AD\) 为两条互相垂直的方向，并设 \(O\) 在这两个方向上的有向坐标分别为 \(u\)，\(v\). 由于 \(A_1O\perp\) 底面，\(AA_1\) 在 \(AB\) 和 \(AD\) 方向上的投影分别等于 \(u\)，\(v\)，因此
\[
u=AA_1\cos\angle A_1AB=3\cos\frac{\pi}{3}=\frac{3}{2},\qquad v=AA_1\cos\angle A_1AD=3\cos\frac{\pi}{3}=\frac{3}{2}.
\]
因为 \(AB\perp AD\)，所以 \(\angle BAD\) 是直角，且其平分线上的点在 \(AB\)，\(AD\) 两个方向上的分量相等. 由 \(u=v=\frac{3}{2}\)，可知 \(O\) 在 \(\angle BAD\) 的平分线上.

 \item 因为 \(AB\perp AD\)，有
\[
AO^2=u^2+v^2=\frac{9}{2}.
\]
又因为三角形 \(AA_1O\) 在 \(O\) 处为直角，所以
\[
A_1O^2=AA_1^2-AO^2=9-\frac{9}{2}=\frac{9}{2},\qquad A_1O=\frac{3}{\sqrt{2}}.
\]
底面 \(ABCD\) 是矩形，面积为 \(AB\cdot AD=5\times4=20\)，故平行六面体的体积为
\[
V=20\times\frac{3}{\sqrt{2}}=30\sqrt{2}.
\]
\end{enumerate}
\end{solution}

\begin{problem}
自点 \(A(-3,3)\) 发出的光线 \(L\) 射到 \(x\) 轴上，被 \(x\) 轴反射，其反射光线所在直线与圆 \(x^2 + y^2 - 4x - 4y + 7 = 0\) 相切，求光线 \(L\) 所在直线的方程.
\end{problem}

\begin{solution}
将点 \(A(-3,3)\) 关于 \(x\) 轴对称，得到点 \(A'(-3,-3)\). 反射光线所在直线关于 \(x\) 轴反射后经过 \(A'\)，且与圆相切. 将圆的方程化为
\[
(x-2)^2+(y-2)^2=1,
\]
可知圆心为 \(C(2,2)\)，半径为 \(1\).

设过 \(A'\) 的切线斜率为 \(m\)，则其方程为 \(y+3=m(x+3)\)，即 \(mx-y+3m-3=0\). 圆心到该直线的距离等于半径，因而
\[
\frac{|2m-2+3m-3|}{\sqrt{m^2+1}}=1.
\]
平方并整理，得 \(25(m-1)^2=m^2+1\)，即 \(12m^2-25m+12=0\)，所以
\[
m=\frac{4}{3}\quad\text{或}\quad m=\frac{3}{4}.
\]
令切线与 \(x\) 轴的交点为 \(Q\)，由 \(y+3=m(x+3)\) 得
\[
Q=\left(\frac{3}{m}-3,0\right).
\]
当 \(m=\frac{4}{3}\) 和 \(m=\frac{3}{4}\) 时，分别有 \(Q=\left(-\frac{3}{4},0\right)\) 和 \(Q=(1,0)\)，二者都在 \(A\) 的右侧，因而入射光线确实可由 \(A\) 射至 \(Q\). 对直线 \(mx-y+3m-3=0\)，圆心 \(C\) 到该直线的垂足为
\[
T=C-\frac{5(m-1)}{m^2+1}(m,-1).
\]
代入两种斜率，得切点分别为 \(T=\left(\frac{6}{5},\frac{13}{5}\right)\) 和 \(T=\left(\frac{13}{5},\frac{6}{5}\right)\)，均位于从相应的 \(Q\) 沿斜率 \(m\) 的正向射线，因此两条候选切线均满足反射光线的方向条件. 关于 \(x\) 轴反射回去后，原光线所在直线经过 \(A(-3,3)\)，斜率分别为 \(-\frac{4}{3}\)，\(-\frac{3}{4}\). 因此
\[
y-3=-\frac{4}{3}(x+3)\quad\text{或}\quad y-3=-\frac{3}{4}(x+3),
\]
即所求直线方程为 \(4x+3y+3=0\) 或 \(3x+4y-3=0\).
\end{solution}

\begin{problem}
已知 \(a > 0\)，\(a \neq 1\)，试求使方程 \(\log_a(x - ak) = \log_{a^2}(x^2 - a^2)\) 有解的 \(k\) 的取值范围.
\end{problem}

\begin{solution}
由对数的定义域，必须有 \(x-ak>0\) 且 \(x^2-a^2>0\). 又因为 \(a>0\)，\(a\neq1\)，所以
\[
\log_{a^2}(x^2-a^2)=\frac{1}{2}\log_a(x^2-a^2)=\log_a\sqrt{x^2-a^2}.
\]
由于 \(\log_a x\) 在其定义域内是一一对应的，原方程等价于
\[
x-ak=\sqrt{x^2-a^2}.
\]
两边均为正数，平方得
\[
(x-ak)^2=x^2-a^2,
\]
从而 \(2akx=a^2(k^2+1)\). 当 \(k=0\) 时上式变为 \(0=a^2\)，不可能；当 \(k\neq0\) 时，解得
\[
x=\frac{a(k^2+1)}{2k},\qquad x-ak=\frac{a(1-k^2)}{2k}.
\]
要使 \(x-ak>0\)，由于 \(a>0\)，必须有 \(\frac{1-k^2}{k}>0\)，即 \(k<-1\) 或 \(0<k<1\). 对这些 \(k\)，由上式可得 \(x^2-a^2=(x-ak)^2>0\)，定义域条件同时满足. 因此 \(k\) 的取值范围为
\[
k\in(-\infty,-1)\cup(0,1).
\]
\end{solution}

\begin{problem}
是否存在常数 \(a\)，\(b\)，\(c\) 使得等式 \(1 \cdot 2^2 + 2 \cdot 3^2 + \cdots + n(n + 1)^2 = \frac{n(n + 1)}{12}(an^2 + bn + c)\) 对一切自然数 \(n\) 都成立？并证明你的结论.
\end{problem}

\begin{solution}
将左端逐项展开，得

\[
\sum_{j=1}^{n}j(j+1)^2=\sum_{j=1}^{n}(j^3+2j^2+j).
\]

由幂和公式，\(\sum_{j=1}^{n}j^3=\left(\frac{n(n+1)}{2}\right)^2\)，\(\sum_{j=1}^{n}j^2=\frac{n(n+1)(2n+1)}{6}\)，\(\sum_{j=1}^{n}j=\frac{n(n+1)}{2}\). 因此
\[
\sum_{j=1}^{n}j(j+1)^2=\left(\frac{n(n+1)}{2}\right)^2+2\cdot\frac{n(n+1)(2n+1)}{6}+\frac{n(n+1)}{2}.
\]
将右侧通分并提取 \(\frac{n(n+1)}{12}\)，得到
\[
\sum_{j=1}^{n}j(j+1)^2=\frac{n(n+1)}{12}\left(3n(n+1)+4(2n+1)+6\right).
\]
将括号内展开，得
\[
3n(n+1)+4(2n+1)+6=3n^2+11n+10.
\]
所以
\[
\sum_{j=1}^{n}j(j+1)^2=\frac{n(n+1)}{12}\left(3n^2+11n+10\right).
\]
所以对一切自然数 \(n\)，题设等式确实成立，取
\[
a=3,\qquad b=11,\qquad c=10.
\]
反过来，若题设等式对一切自然数 \(n\) 成立，则对一切正整数 \(n\) 也成立. 由于正整数 \(n\) 满足 \(n(n+1)\neq0\)，约去 \(\frac{n(n+1)}{12}\) 后，必须有
\[
an^2+bn+c=3n^2+11n+10
\]
对一切正整数 \(n\) 成立，即
\[
(a-3)n^2+(b-11)n+(c-10)=0
\]
对一切正整数 \(n\) 成立. 这个次数不超过 \(2\) 的多项式有无穷多个根，只能恒等于零. 于是
\[
a=3,\qquad b=11,\qquad c=10.
\]
所以存在这样的常数，且唯一为 \(a=3\)，\(b=11\)，\(c=10\).
\end{solution}

\begin{problem}
设 \(f(x)\) 是定义在区间 \((-\infty, +\infty)\) 上以 \(2\) 为周期的函数，对 \(k \in \Z\)，用 \(I_k\) 表示区间 \((2k - 1, 2k + 1]\)，已知当 \(x \in I_0\) 时，\(f(x) = x^2\).
\begin{enumerate}
\item 求 \(f(x)\) 在 \(I_k\) 上的解析表达式；
\item 对自然数 \(k\)，求集合 \(M_k = \left\{ a \mid \text{使方程 } f(x) = ax \text{ 在 } I_k \text{ 上有两个不等的实根} \right\}\).
\end{enumerate}
\end{problem}

\begin{solution}
\begin{enumerate}
\item 当 \(x\in I_k\) 时，\(x-2k\in(-1,1]=I_0\). 由 \(f(x)\) 的周期为 \(2\)，有 \(f(x)=f(x-2k)\)，再由 \(f(t)=t^2\)（\(t\in I_0\)）可得
\[
f(x)=(x-2k)^2,\qquad x\in I_k.
\]

\item 按题目中自然数 \(k=1,2,\ldots\) 讨论. 令 \(t=x-2k\)，则 \(t\in(-1,1]\)，且 \(x=2k+t>0\). 因而当 \(a<0\) 时方程没有实根；当 \(a=0\) 时方程化为 \(t^2=0\)，只有一个实根. 对 \(a>0\)，方程 \(f(x)=ax\) 等价于
\[
t^2=a(2k+t),\qquad\text{即}\qquad a=h(t):=\frac{t^2}{2k+t}.
\]
在 \((-1,1]\) 上，
\[
h'(t)=\frac{t(4k+t)}{(2k+t)^2}.
\]
因为 \(k\geqslant1\)，所以 \(4k+t>0\)，从而 \(h(t)\) 在 \((-1,0)\) 上递减，在 \((0,1]\) 上递增. 并且
\[
\lim_{t\to-1^+}h(t)=\frac{1}{2k-1},\qquad h(0)=0,\qquad h(1)=\frac{1}{2k+1}.
\]
因此，负半段给出值域 \(\left(0,\frac{1}{2k-1}\right)\)，正半段给出值域 \(\left(0,\frac{1}{2k+1}\right]\). 方程有两个不等的实根，当且仅当同一个正数 \(a\) 同时落在这两个值域中，即
\[
0<a\leqslant\frac{1}{2k+1}.
\]
所以
\[
M_k=\left(0,\frac{1}{2k+1}\right].
\]
若将 \(0\) 也列入自然数，则另由 \(I_0=(-1,1]\) 直接得到 \(M_0=(-1,0)\cup(0,1]\).
\end{enumerate}
\end{solution}
